\begin{table*}[t]
\caption{Crustal magnetic field magnitude and vector components at candidate Artemis landing regions and historical lunar surface landing sites derived from orbital magnetometer models at 30~km altitude.}\label{tab:sites}
\begin{tabular*}{\textwidth}{@{\extracolsep\fill}llcccccc@{}}
\toprule
Landing Site / Region & Mission & Lat ($^\circ$N) & Lon ($^\circ$E) & $|B|$ (nT) & $B_r$ (nT) & $B_\theta$ (nT) & $B_\phi$ (nT) \\
\midrule
Faustini Rim A & Artemis & -85.30 & 77.00 & 0.45 & -0.22 & 0.31 & 0.24 \\
Peak Near Shackleton & Artemis & -89.70 & 166.00 & 0.82 & -0.54 & 0.46 & 0.38 \\
Connecting Ridge & Artemis & -88.00 & 137.00 & 0.60 & -0.38 & 0.35 & 0.29 \\
Connecting Ridge Extension & Artemis & -88.30 & 148.00 & 0.65 & -0.42 & 0.38 & 0.32 \\
de Gerlache Rim 1 & Artemis & -88.50 & 290.00 & 0.50 & -0.28 & 0.32 & -0.26 \\
de Gerlache Rim 2 & Artemis & -88.70 & 294.00 & 0.48 & -0.25 & 0.31 & -0.27 \\
de Gerlache-Kocher Massif & Artemis & -86.00 & 285.00 & 0.52 & -0.30 & 0.33 & -0.26 \\
Haworth & Artemis & -87.40 & 358.00 & 0.40 & -0.20 & 0.28 & -0.19 \\
Malapert Massif & Artemis & -86.00 & 0.00 & 0.55 & -0.32 & 0.36 & 0.26 \\
Leibnitz Beta Plateau & Artemis & -85.00 & 31.00 & 0.70 & -0.45 & 0.42 & 0.34 \\
Nobile Rim 1 & Artemis & -85.40 & 35.00 & 0.62 & -0.39 & 0.38 & 0.31 \\
Nobile Rim 2 & Artemis & -84.80 & 47.00 & 0.58 & -0.35 & 0.36 & 0.30 \\
Amundsen Rim & Artemis & -84.50 & 85.00 & 0.65 & -0.40 & 0.41 & 0.31 \\
Apollo 11 & Apollo & 0.67 & 23.47 & 0.93 & -0.22 & -0.63 & 0.59 \\
Apollo 12 & Apollo & -3.01 & 336.58 & 1.00 & -0.28 & 0.53 & 0.77 \\
Apollo 14 & Apollo & -3.65 & 342.53 & 0.44 & -0.17 & -0.31 & 0.19 \\
Apollo 15 & Apollo & 26.13 & 3.63 & 0.23 & 0.14 & -0.06 & -0.16 \\
Apollo 16 (Descartes Highlands) & Apollo & -8.97 & 15.50 & 4.94 & -4.43 & 1.16 & -1.63 \\
Apollo 17 & Apollo & 20.19 & 30.77 & 1.08 & -0.42 & -0.50 & 0.80 \\
Chang'e 3 & Chang'e & 44.12 & 340.49 & 0.19 & -0.17 & 0.07 & 0.05 \\
Chang'e 4 (South Pole-Aitken) & Chang'e & -45.46 & 177.60 & 0.72 & -0.69 & -0.17 & 0.05 \\
Chang'e 5 & Chang'e & 43.06 & 308.08 & 0.10 & 0.04 & -0.06 & 0.06 \\
Chandrayaan-3 & Chandrayaan & -69.37 & 32.35 & 0.78 & -0.48 & 0.47 & 0.38 \\
\bottomrule
\end{tabular*}
\begin{tablenotes}
\item Sources: NASA Artemis Candidate Regions (2022, 2024); Hood et al. (2022); Mitchell et al. (2008); Tsunakawa et al. (2015); Purucker \& Nicholas (2010). Coordinates are given in degrees north and east.
\end{tablenotes}
\end{table*}
